\section*{Chapitre \ref{ch:pl:the-platform-map} --- La carte de la plateforme}

\begin{solution}{exo:pl:the-platform-map:1}
Il démarre quand sa dernière entrée est prête~: les positions, à 20~h~15. Il se termine à 20~h~40. Les valorisations (17~h~05) et le référentiel des titres (security master)
(19~h~55) étaient prêts plus tôt~; ce sont les positions qui décident.
\end{solution}

\begin{solution}{exo:pl:the-platform-map:2}
La date limite d'arrivée sûre du fichier est 22~h~35~; à 22~h~50, il a 15 minutes de retard. La copie de référence (golden copy) s'exécute de 22~h~50 à 23~h~15, les variables
jusqu'à 01~h~15, et les backtests se terminent à 07~h~15~: le rapport de recherche accuse 15 minutes de retard. Tout autre rapport conserve encore une marge (le
batch de risque, par exemple, s'étend de 23~h~35 à 05~h~35).
\end{solution}

\begin{solution}{exo:pl:the-platform-map:3}
Seuls deux jeux de données en dérivent~: les écarts de rapprochement et le P\&L signé, car il n'est lu que par le rapprochement.
Un petit \omterm{def:pl:the-platform-map:blast}{rayon d'impact} (blast radius) indique jusqu'où une erreur se propage dans le batch, et non son ampleur~: un relevé de compensation erroné peut
masquer un véritable écart dans les positions de la société, c'est précisément ce que le rapprochement sert à détecter. Il faut le vérifier au
point d'utilisation (chapitre~22), alors que peu de choses reposent dessus.
\end{solution}

\begin{solution}{exo:pl:the-platform-map:4}
Les positions sont prêtes 45 minutes après le fichier (25 minutes de copie de référence, 20 de positions), si bien que le batch
ne lit des entrées complètes que si le fichier arrive au plus tard à $270 - 45 = 225$ minutes après la clôture, soit 19~h~45. Le fichier nominal arrive à 19~h~30~: 15 minutes
de marge au lieu de 45. Un démarrage plus tôt rend la lecture de données périmées trois fois plus probable lors d'une nuit un peu tardive.
\end{solution}

\begin{solution}{exo:pl:the-platform-map:5}
Les valorisations vérifiées sont calculées à partir d'autres entrées (cotations des courtiers, données de consensus) par un autre système et comportent des ajustements~; elles
ne sont pas dérivées par la carte des valorisations de la bibliothèque, si bien qu'il s'agit d'un autre jeu de données, avec son propre \omterm{def:pl:the-platform-map:sor}{système de référence} (system of record), le product
control. La validation lit les valorisations vérifiées et rapporte l'écart entre les deux sous forme de nombre, car le but de
la vérification indépendante est que les valorisations propres au desk ne bouclent pas les comptes.
\end{solution}

\begin{solution}{exo:pl:the-platform-map:6}
(a) Informations tardives~: les annulations, les corrections et les exécutions tardives arrivent après que le \omterm{def:pl:the-platform-map:batch}{traitement en flux} a appliqué les originaux, et un drop copy peut
montrer des exécutions que la session de trading a manquées (chapitre~17). (b) Données de référence~: le batch applique les opérations sur titres du soir, issues du
référentiel des titres (un split modifie les quantités), pas le \omterm{def:pl:the-platform-map:batch}{traitement en flux}. Les livres sont clôturés sur le batch, rapprochés avec le
relevé de compensation~; l'écart avec le \omterm{def:pl:the-platform-map:batch}{traitement en flux} est traité comme un écart de rapprochement.
\end{solution}

\begin{solution}{exo:pl:the-platform-map:7}
Son dernier démarrage possible est à $840 - 60 = 780$ minutes après la clôture, soit 05~h~00. Le dernier démarrage possible d'aucun autre job ne se déplace~: les positions devaient
déjà être prêtes à 01~h~00 pour le batch de risque, bien avant 05~h~00. La date limite d'arrivée sûre du relevé de compensation passe à 05~h~00 au lieu de
05~h~50, car la prévision le lit désormais avec une échéance plus proche que celle du rapprochement.
\end{solution}

\begin{solution}{exo:pl:the-platform-map:8}
Le batch de risque ne figure pas sur la chaîne la plus longue. Ramené à trois heures, il se termine à 23~h~15 au lieu de 02~h~15, mais le batch se termine toujours à
03~h~55, avec les backtests~: la chaîne la plus longue est celle du fichier des opérations sur titres, de la copie de référence, des variables et des backtests. Et
doubler le nombre de machines ne réduit de moitié un job que s'il se parallélise parfaitement et que rien d'autre ne le limite (chapitres~13 et~20).
\end{solution}

\begin{solution}{pb:pl:the-platform-map:1}
\begin{enumerate}
\item Copie de référence 19~h~30--19~h~55~; positions 19~h~55--20~h~15~; batch de risque 20~h~15--02~h~15.
\item Les backtests, à 03~h~55.
\item Le fichier des opérations sur titres, la copie de référence, les variables et les backtests~: $25 + 120 + 360$ minutes à partir de 19~h~30. La chaîne propre au
batch de risque (fichier, copie de référence, positions, batch de risque) se termine à 02~h~15, 100 minutes plus tôt, car la branche de recherche
porte deux jobs longs.
\item $900 - 715 = 185$ minutes.
\item 285 minutes~: son dernier démarrage possible est à 01~h~00 et il démarre à 20~h~15.
\item 22~h~35, fixé par la copie de référence, dont le dernier démarrage possible est fixé par les variables (dernier démarrage 23~h~00) et les backtests (01~h~00).
\item 22~h~25~: la saisie des transactions alimente la compaction des ticks, qui alimente les variables et donc les backtests~; les exécutions alimentent les positions,
dont le dernier démarrage possible (00~h~40) est fixé par le batch de risque, une échéance plus tardive dans la chaîne.
\item Les exécutions~: 515 minutes (de 16~h~05 à 00~h~40).
\item Neuf~: le référentiel des titres, les positions, le P\&L, le rapport de risque, les écarts de rapprochement, le P\&L signé, les
rapports de transaction, les variables et les backtests.
\item Les écarts de rapprochement et le P\&L signé.
\item 21~h~55 et 22~h~15.
\item Le référentiel des titres et les positions (les scénarios et les valorisations étaient prêts bien avant).
\item De 22~h~15 à 22~h~40, sur le nouveau référentiel des titres et les nouvelles positions.
\item 20~h~15~: le fichier doit arriver 45 minutes avant le démarrage à 21~h~00.
\item De 22~h~15 à 04~h~15, soit 165 minutes avant l'échéance.
\item \textbf{Résultat nommé.} Une nuit nominale, le batch se termine à 03~h~55, 185 minutes avant 07~h~00, sur la chaîne de recherche~; le
fichier des opérations sur titres peut arriver jusqu'à 22~h~35~; un batch de risque lancé à heure fixe à 21~h~00 lit un référentiel des titres périmé et des
positions périmées dès que le fichier arrive après 20~h~15, soit 45 minutes après son heure nominale, alors que lancé sur ses entrées il aurait
terminé à 04~h~15 la nuit de l'accroche.
\item Les versions périmées existent et sont valides~; un job à démarrage fixe n'a aucun moyen de savoir qu'une version plus récente arrive.
\item La troisième cause~: deux versions des mêmes jeux de données, lues à des moments différents.
\item Faire vérifier par le job que la version de chaque entrée pour la date de gestion est complète avant son démarrage (un marqueur d'achèvement
écrit par le producteur), et refuser ou attendre sinon~; ou le laisser démarrer mais enregistrer la version de chaque entrée lue, de sorte que
le rapport soit marqué comme périmé et relancé à l'arrivée de la nouvelle version.
\item La carte peut être vérifiée et simulée~: un programme indique ce qu'un fichier atteint, quand chaque rapport est prêt, et quel job lit une
entrée incomplète.
\end{enumerate}
\end{solution}

\begin{solution}{iq:pl:the-platform-map:1}
Le seul système dont la copie d'un jeu de données fait autorité~: les corrections y sont faites, et tout le reste en est dérivé.
Sans un système par jeu de données, deux systèmes peuvent chacun avoir raison selon leurs propres règles et se contredire, et personne ne peut dire quel rapport est
juste.
\iqlookfor{Corrections en un seul lieu~; copies dérivées et non modifiées~; le désaccord comme symptôme de deux responsables.}
\end{solution}

\begin{solution}{iq:pl:the-platform-map:2}
En flux lorsqu'une décision en a besoin maintenant (risque pré-trade, écran du trader, kill switch)~; en batch lorsque la réponse doit être
complète et définitive (les livres, le rapprochement de compensation, les rapports réglementaires), car les exécutions tardives, les annulations et les opérations sur titres
ne sont connues qu'après la clôture. La plupart des sociétés font les deux et les rapprochent.
\iqlookfor{Latence contre complétude~; les deux conservés et rapprochés.}
\end{solution}

\begin{solution}{iq:pl:the-platform-map:3}
Remonter les deux à travers la carte jusqu'aux jeux de données qu'elles lisent, ainsi que les versions et les moments de lecture. Vérifier les trois causes dans
l'ordre~: deux \omterm{def:pl:the-platform-map:sor}{systèmes de référence} pour une même entrée, une copie faite hors des flux déclarés (un tableur, un cache), ou le même jeu de données
lu à deux moments différents (un job exécuté avant qu'une entrée tardive soit complète). Corriger ensuite la cause, et non le chiffre.
\iqlookfor{Traçabilité et versions plutôt que supposition~; la cause temporelle.}
\end{solution}

\begin{solution}{iq:pl:the-platform-map:4}
Sur ses entrées, lorsque c'est possible~: il démarre le plus tôt possible et ne lit jamais une entrée incomplète, mais il a besoin d'un
signal d'achèvement fiable de chaque producteur, et il peut démarrer à un moment mal choisi. Un démarrage fixe est prévisible pour les machines et les personnes,
mais il lit silencieusement des données périmées lorsqu'une entrée est en retard. Si un démarrage fixe est conservé, vérifier l'achèvement avant le lancement et enregistrer
les versions lues.
\iqlookfor{La défaillance par lecture de données périmées nommée~; les marqueurs d'achèvement.}
\end{solution}

\begin{solution}{iq:pl:the-platform-map:5}
Un parcours à rebours du graphe de dépendances~: à partir de l'échéance de chaque rapport, on retranche les durées des jobs le long de chaque chemin jusqu'au
fichier~; le plus tôt des derniers démarrages possibles qui en résultent, parmi les jobs qui le lisent, donne la date limite d'arrivée sûre du fichier. Mesurer la
dispersion des durées et conserver une marge, car le parcours suppose des durées fixes et des machines libres.
\iqlookfor{Dernier démarrage possible, en remontant le graphe, avec une marge pour les durées variables.}
\end{solution}

\begin{solution}{iq:pl:the-platform-map:6}
Partir des sorties sur lesquelles on s'appuie (le P\&L, le rapport de risque, les rapports réglementaires) et remonter ce que chacune lit, depuis
où et quand, en le consignant sous forme de données plutôt que sous forme de schéma~; faire déclarer leurs entrées par les systèmes pour que la carte soit générée et
vérifiée, et non maintenue à la main. Premiers contrôles~: un \omterm{def:pl:the-platform-map:sor}{système de référence} par jeu de données, aucune copie non déclarée, et des entrées complètes
pour chaque job planifié au moment de son démarrage.
\iqlookfor{Carte sous forme de données, dérivée des systèmes~; d'abord les contrôles de propriété et de temporalité.}
\end{solution}
